\section{Results}

The proposed hybrid retrieval and reranking framework was evaluated using 25 natural-language queries related to biomedical natural language processing, transformer-based clinical language models, electronic health records, ICD code prediction, BioBERT, ClinicalBERT, and large language models in healthcare. For each query, the retrieval stage returned 20 candidate evidence chunks, resulting in 500 reranked evidence chunks across the complete evaluation set. During answer generation, the top five reranked chunks were used as the evidence context for each query. The generated answers were subsequently evaluated using claim-level grounding assessment. This result should be interpreted within the scope of the evaluated 25-query test set and the automated judge-based evaluation protocol.

Table~\ref{tab:overall_grounding_summary} summarizes the overall evaluation results. Across the 25 generated answers, the grounding evaluation process extracted 200 factual claims. The number of claims per query ranged from 1 to 12, with an average of 8.0 claims per query. All 200 claims were judged to be supported by the retrieved evidence, while no claims were marked as unsupported. Therefore, the overall claim-level grounding accuracy was 100.0\%. At the query level, all 25 queries achieved 100.0\% grounding accuracy, indicating that every factual claim extracted from each generated response was supported by the corresponding retrieved evidence.

\begin{table}[!t]
\centering
\caption{Overall claim-level grounding evaluation results.}
\label{tab:overall_grounding_summary}
\begin{tabular}{lc}
\hline
\textbf{Evaluation metric} & \textbf{Value} \\
\hline
Total evaluated queries & 25 \\
Total reranked evidence chunks & 500 \\
Reranked chunks per query & 20 \\
Chunks used for answer generation per query & 5 \\
Total extracted factual claims & 200 \\
Mean claims per query & 8.0 \\
Claim-count range per query & 1--12 \\
Supported claims & 200 \\
Unsupported claims & 0 \\
Overall grounding accuracy & 100.0\% \\
Queries with 100\% grounding accuracy & 25 of 25 \\
\hline
\end{tabular}
\end{table}

The evidence retrieved by the system was distributed across five source documents, as shown in Table~\ref{tab:source_distribution}. The largest proportion of reranked chunks came from the review paper on transformers and large language models in healthcare, which contributed 155 chunks, corresponding to 31.0\% of the retrieved evidence. The second-largest contribution came from the transformer models in biomedicine paper, which contributed 123 chunks, or 24.6\%. The lightweight transformers for clinical natural language processing paper contributed 87 chunks, or 17.4\%. The ICD prediction using ClinicalBERT paper contributed 70 chunks, or 14.0\%, while the BioBERT paper contributed 65 chunks, or 13.0\%. This distribution shows that the retrieval pipeline used evidence from multiple documents rather than relying on a single dominant source.

\begin{table}[!t]
\centering
\caption{Distribution of reranked evidence chunks across source documents.}
\label{tab:source_distribution}
\begin{tabular}{lcc}
\hline
\textbf{Source document} & \textbf{Chunks} & \textbf{Percentage} \\
\hline
Transformers and LLMs in healthcare review & 155 & 31.0\% \\
Transformer models in biomedicine & 123 & 24.6\% \\
Lightweight transformers for clinical NLP & 87 & 17.4\% \\
ICD prediction using ClinicalBERT & 70 & 14.0\% \\
BioBERT for biomedical text mining & 65 & 13.0\% \\
\hline
\textbf{Total} & \textbf{500} & \textbf{100.0\%} \\
\hline
\end{tabular}
\end{table}

The reranking stage substantially changed the initial retrieval order. Among the 500 retrieved chunks, only 44 chunks retained the same rank after reranking, corresponding to 8.8\% of the evidence set. For the top-ranked reranked chunks, the original Bedrock rank varied from 1 to 16, with a mean original rank of 6.0. Only 10 of the 25 final top-ranked chunks were originally ranked within the top three Bedrock retrieval results, while three final top-ranked chunks were originally ranked outside the top ten. These findings indicate that the reranking stage had a meaningful effect on evidence prioritization by reordering initially retrieved chunks before answer generation.

\begin{table}[!t]
\centering
\caption{Effect of reranking on retrieved evidence ordering.}
\label{tab:reranking_effect}
\begin{tabular}{lc}
\hline
\textbf{Reranking metric} & \textbf{Value} \\
\hline
Total reranked chunks & 500 \\
Chunks with unchanged rank & 44 \\
Percentage of unchanged ranks & 8.8\% \\
Mean original Bedrock rank of final top-ranked chunks & 6.0 \\
Minimum original Bedrock rank of final top-ranked chunks & 1 \\
Maximum original Bedrock rank of final top-ranked chunks & 16 \\
Final top-ranked chunks originally in Bedrock top 3 & 10 of 25 \\
Final top-ranked chunks originally outside Bedrock top 10 & 3 of 25 \\
\hline
\end{tabular}
\end{table}

%The initial Bedrock retrieval score across all 500 chunks had a mean value of 0.501 with a standard deviation of 0.059. The Cohere rerank score had a mean value of 0.442 with a standard deviation of 0.242. Since these two scores are generated by different components, they should not be interpreted as directly comparable on the same absolute scale. However, within the reranking output, the final top-ranked chunks received relatively high rerank scores, with a mean score of 0.801 and a range from 0.566 to 0.919. This suggests that the reranker assigned stronger relevance priority to the chunks selected at the top of the final evidence list.

The initial Bedrock retrieval score across all 500 retrieved chunks had a mean value of 0.501 with a standard deviation of 0.059. The Cohere rerank score across the same 500 reranked chunks had a mean value of 0.442 with a standard deviation of 0.242. Since these two scores are generated by different components, they should not be interpreted as directly comparable on the same absolute scale. However, within the Cohere reranking output, the 25 final top-ranked chunks, one for each query, received relatively high rerank scores, with a mean score of 0.801, a standard deviation of 0.098, and a range from 0.566 to 0.919. This suggests that the reranker assigned stronger relevance priority to the chunks selected at the top of the final evidence list.

\begin{table}[!t]
\centering
\caption{Retrieval and reranking score summary.}
\label{tab:score_summary}
\begin{tabular}{lcccc}
\hline
\textbf{Score type} & \textbf{Mean} & \textbf{Std.} & \textbf{Min.} & \textbf{Max.} \\
\hline
Bedrock retrieval score & 0.501 & 0.059 & 0.401 & 0.727 \\
Cohere rerank score & 0.442 & 0.242 & 0.016 & 0.919 \\
Top-ranked rerank score & 0.801 & .098 & 0.566 & 0.919 \\
\hline
\end{tabular}
\end{table}

The claim-level grounding results demonstrate that the generated answers were strongly aligned with the retrieved evidence. Because each generated answer was decomposed into individual factual claims, the evaluation provides a more fine-grained assessment than answer-level correctness alone. A generated response may contain both supported and unsupported statements; therefore, claim-level evaluation provides a stricter way to examine whether the response remains grounded in the retrieved evidence. In this experiment, all extracted claims were supported, indicating that the combination of hybrid retrieval, reranking, conservative prompting, and independent judge-model evaluation produced highly grounded responses for the selected biomedical NLP query set.

It is important to interpret these results as internal grounding accuracy rather than strict citation accuracy. The generated answers did not include explicit source labels in the final response. Therefore, the reported 100.0\% accuracy indicates that each factual claim was supported by at least one retrieved evidence chunk associated with the query. Strict citation accuracy would require each generated claim to include an explicit citation label and would further evaluate whether the cited source specifically supports that claim. Thus, the current experiment demonstrates strong evidence grounding, while future work should extend the framework to evaluate explicit sentence-level citation correctness.

Overall, the results show that the proposed hybrid retrieval and reranking framework produced evidence-grounded answers across all evaluated biomedical NLP and healthcare transformer queries. The reranking stage meaningfully changed the initial retrieval order, the retrieved evidence was distributed across multiple source documents, and the claim-level evaluation confirmed full support for all generated factual claims. These findings suggest that the framework can support reliable retrieval-augmented answer generation when the source document collection contains sufficient evidence for the target query. Future evaluation should include a larger and more diverse query set, human expert validation, explicit citation-label generation, and comparison against retrieval-only baselines to determine whether reranking improves both grounding accuracy and strict citation accuracy.



\subsection{Estimated Computational Time}

\begin{table}[!t]
\centering
\caption{Estimated execution time in Google Colab with NVIDIA A100 GPU.}
\label{tab:estimated_runtime}
\begin{tabular}{lc}
\hline
\textbf{Pipeline stage} & \textbf{Estimated time} \\
\hline
Environment setup and package loading & 1--3 min \\
Hybrid retrieval for 25 queries & 2--5 min \\
Cohere reranking of retrieved chunks & 3--6 min \\
Grounded answer generation & 5--8 min \\
Claim extraction and grounding evaluation & 8--12 min \\
CSV export and summary statistics & $<$1 min \\
\hline
\textbf{Estimated total runtime} & \textbf{20--35 min} \\
\hline
\end{tabular}
\end{table}

The experimental pipeline was executed in a Google Colab environment with an NVIDIA A100 GPU. For the evaluated setting of 25 queries, 500 reranked evidence chunks, and 200 claim-level grounding judgments, the estimated end-to-end wall-clock execution time was approximately 20--35 minutes. This estimate includes hybrid retrieval, reranking, answer generation, claim extraction, grounding evaluation, and result-file generation. It should be noted that the Amazon Bedrock and Cohere model calls were executed through external managed services; therefore, the total runtime was influenced primarily by API latency, network response time, model availability, and request processing limits rather than by GPU computation alone. The A100 GPU mainly supported the local notebook execution, data handling, and result processing steps. Although the experiment was conducted in a Google Colab A100 GPU environment, the main computational bottleneck was external model inference through Amazon Bedrock and Cohere APIs rather than local GPU computation.



